\section{Largest Divisible Subset}
\label{sec:dp-largest-divisible-subset}

\problemstatement{Given an array of \textbf{distinct} positive
integers, find the largest subset such that for every pair of elements
$(a,b)$ in it, either $a \bmod b = 0$ or $b \bmod a = 0$.}

\begin{patternbox}
\emph{``Largest subset where every pair divides''} is
Section~\ref{sec:dp-lis-intro}'s ending-here template with the
condition $\textit{arr}[j] < \textit{arr}[i]$ replaced by
$\textit{arr}[i] \bmod \textit{arr}[j] = 0$ -- after first
\textbf{sorting} the array, exactly the same preparatory step LIS uses
implicitly (LIS's input is typically already read left to right; here
we must sort first to make the pairwise-divisibility condition reduce
to a simple left-to-right chain check).
\end{patternbox}

\subsection*{Why sorting first is essential}

After sorting, if $\textit{arr}[k]$ divides $\textit{arr}[j]$ and
$\textit{arr}[j]$ divides $\textit{arr}[i]$ (with $k<j<i$ in the sorted
order), then $\textit{arr}[k]$ \emph{automatically} divides
$\textit{arr}[i]$ too (divisibility is transitive). So checking only
\textbf{consecutive} chain links ($\textit{arr}[j] \mid
\textit{arr}[i]$ for the immediately preceding chosen element) is
sufficient to guarantee \emph{every pair} in the resulting chain
satisfies the divisibility condition -- exactly mirroring why LIS only
checks $\textit{arr}[j]<\textit{arr}[i]$ pairwise rather than every
triple.

\subsection*{The recurrence}

\[
  \textit{dp}[i] = 1 + \max_{j<i,\ \textit{arr}[i] \bmod
  \textit{arr}[j] = 0} \textit{dp}[j] \quad (\text{or } 1 \text{ if no
  such } j).
\]

\subsection*{Algorithm (with reconstruction)}

Exactly Section~\ref{sec:dp-print-lis}'s parent-pointer approach,
with the extension condition swapped to divisibility.

\begin{lstlisting}
#include <bits/stdc++.h>
using namespace std;

vector<int> largestDivisibleSubset(vector<int> arr) {
    int n = arr.size();
    sort(arr.begin(), arr.end());

    vector<int> dp(n, 1), parent(n);
    for (int i = 0; i < n; i++) parent[i] = i;

    int lastIndex = 0;
    for (int i = 0; i < n; i++) {
        for (int j = 0; j < i; j++) {
            if (arr[i] % arr[j] == 0 && 1 + dp[j] > dp[i]) {
                dp[i] = 1 + dp[j];
                parent[i] = j;
            }
        }
        if (dp[i] > dp[lastIndex]) lastIndex = i;
    }

    vector<int> result;
    int cur = lastIndex;
    while (parent[cur] != cur) {
        result.push_back(arr[cur]);
        cur = parent[cur];
    }
    result.push_back(arr[cur]);
    reverse(result.begin(), result.end());
    return result;
}

int main() {
    vector<int> arr = {1, 16, 7, 8, 4};
    for (int x : largestDivisibleSubset(arr)) cout << x << " ";
    cout << endl; // 1 4 8 16
    return 0;
}
\end{lstlisting}

\subsection*{Dry run}

\dryrun{Take $\textit{arr} = [1,16,7,8,4]$, sorted: $[1,4,7,8,16]$.}

$\textit{dp}=[1,2,1,3,4]$: $4 \bmod 1=0 \Rightarrow \textit{dp}[1]=2$;
$8 \bmod 4=0 \Rightarrow \textit{dp}[3]=3$ (extending the chain
$1,4,8$); $16 \bmod 8=0 \Rightarrow \textit{dp}[4]=4$ (extending
$1,4,8,16$). Note $7$ extends nothing (no earlier sorted element
divides it), so $\textit{dp}[2]=1$. The maximum is $4$, giving the
subset $[1,4,8,16]$.

\begin{complexitybox}
\textbf{Time Complexity.} $O(n \log n)$ to sort, plus $O(n^2)$ to fill
the table: $O(n^2)$ overall.
\textbf{Space Complexity.} $O(n)$ for \textit{dp} and \textit{parent}.
\end{complexitybox}

\begin{takeawaybox}
Whenever a pairwise condition is \textbf{transitive} (if $A$ relates to
$B$ and $B$ relates to $C$ then $A$ relates to $C$), sorting first and
then checking only consecutive chain links is enough -- the exact
property that lets both LIS's ``less than'' and this problem's
``divides'' collapse an all-pairs condition down to a simple
left-to-right ending-here DP.
\end{takeawaybox}
